\documentclass[11pt,letterpaper]{article}
\usepackage[utf8]{inputenc}
\usepackage[T1]{fontenc}
\usepackage[spanish,es-nodecimaldot]{babel}
\usepackage[margin=1.75cm]{geometry}
\usepackage{graphicx}
\usepackage{booktabs}
\usepackage{tabularx}
\usepackage{amsmath}
\usepackage{amssymb}
\usepackage{float}
\usepackage{caption}
\usepackage{microtype}
\usepackage{xcolor}
\usepackage{hyperref}

\hypersetup{colorlinks=true,linkcolor=blue!45!black,urlcolor=blue!45!black}
\captionsetup{font=small,skip=4pt}
\setlength{\parindent}{0pt}
\setlength{\parskip}{0.28em}
\setlength{\textfloatsep}{8pt}
\setlength{\floatsep}{7pt}
\setlength{\emergencystretch}{2em}
\renewcommand{\arraystretch}{1.08}

\input{report_assets/resumen/resultados.tex}

\title{\vspace{-1.4cm}\textbf{Cadena de suministro de ayuda humanitaria}\\
\large Modelo y resultados del Grupo 3}
\author{Grupo 3}
\date{}

\begin{document}
\maketitle
\vspace{-1.2cm}

\section{Justificación del paradigma}

Se eligió una \textbf{simulación de eventos discretos con estado actualizado en
pasos de seis horas}. La decisión se basó en cuatro criterios: el sistema cambia
por eventos fechados (consumo, reposición y llegada potencial de convoyes), los
inventarios deben conservarse entre bloques, las restricciones de rutas y flota
son operativas, y las conclusiones requieren comparar realizaciones
estocásticas. La dinámica de sistemas sería útil para estudiar flujos agregados,
pero ocultaría la identidad de centros, rutas y vehículos. Un modelo basado en
agentes permitiría conductas individuales, aunque el Excel no proporciona reglas
de decisión para personas o conductores que justifiquen esa complejidad.

La unidad de evolución es el bloque temporal. En cada realización se actualizan
inventario, demanda satisfecha y déficit para cada zona--suministro. La red y la
flota se conservan como entidades explícitas para evaluar factibilidad del
convoy, combustible y distribución. La frontera híbrida es deliberada: el
balance de inventario es discreto y estocástico, mientras la recomendación de
rutas es una heurística posterior que ordena necesidad y eficiencia. No se
presenta esa heurística como un óptimo matemático.

\section{Descripción del modelo (ODD simplificado)}

\textbf{Propósito y alcance.} El modelo estima, durante 72 horas post-evento, el
balance de agua, alimentos y medicamentos, el primer agotamiento y el margen de
respuesta logística. Se ejecutan \NumeroRealizaciones{} realizaciones con una
semilla base reproducible y semillas independientes por realización.

\textbf{Entidades y atributos.} Los cuatro centros de acopio tienen zona,
capacidad e inventario por suministro. Las cinco zonas reciben una demanda por
bloque. Cada suministro posee tasa per cápita y unidad. Las nueve rutas tienen
origen, destino, distancia, tiempo, capacidad vehicular y estado. La flota tiene
tipo, cantidad, capacidad, rendimiento y combustible. Un convoy es una
alternativa evaluada con ruta, vehículo, cantidad, salida y llegada.

\textbf{Variables de estado.} Para zona $z$, suministro $s$ y bloque $t$ se
registran inventario inicial $I^{\mathrm{ini}}_{zst}$, reposición $R_{zst}$,
demanda $D_{zst}$, consumo satisfecho $C_{zst}$, inventario final
$I^{\mathrm{fin}}_{zst}$, déficit $NS_{zst}$ y superávit $S_{zst}$. La demanda
se obtiene con la población $P_{zt}$ y la tasa de consumo $c_s$:
\begin{align}
D_{zst}&=P_{zt}c_s,\\
A_{zst}&=I^{\mathrm{ini}}_{zst}+R_{zst},\\
C_{zst}&=\min(A_{zst},D_{zst}),\\
I^{\mathrm{fin}}_{zst}&=A_{zst}-C_{zst},\\
NS_{zst}&=\max(0,D_{zst}-A_{zst}),\qquad
S_{zst}=\max(0,A_{zst}-D_{zst}).
\end{align}
El agotamiento crítico exige simultáneamente $I^{\mathrm{fin}}_{zst}=0$ y
$NS_{zst}>0$; así no se confunde inventario exactamente nulo con escasez.

\textbf{Logística y combustible.} Para una ruta de distancia de ida $d_r$, un
vehículo con consumo $g_v$ y $n$ unidades movilizadas usa
\begin{equation}
G_{rv}=2d_r g_v n.
\end{equation}
La frecuencia recomendada se limita por viajes completos dentro del bloque y
por la reserva de combustible compartida. Una ruta no disponible no participa;
si el Grupo 4 recomienda vehículo liviano, el camión grande queda excluido.

\textbf{Scheduling, comunicación e inicialización.} En cada bloque se recibe la
reposición programada, se calcula la demanda, se atiende hasta agotar la
disponibilidad y el inventario final pasa al siguiente bloque. Los centros no
intercambian mensajes como agentes: la ``comunicación'' es el estado tabular de
demanda, rutas y llegada de reposiciones. Inventarios, rutas, flota, combustible
y reposiciones provienen del Excel; la población formal post-intercambio viene
del Grupo 1. El Grupo 4 actualiza tiempos desde Z4 y compatibilidad de vehículo,
sin crear capacidades ni rutas ausentes.

\section{Resultados y análisis}

\subsection{Trayectorias, balance y agotamiento}

La Figura~\ref{fig:trayectorias} muestra la media y el IC 95\% del inventario en
las combinaciones críticas. El escenario base identifica
\ZonaCriticaBase{}/medicamentos; después del intercambio la combinación crítica
abastecible es \ZonaCriticaPost{}/\SuministroCritico{}. Su agotamiento medio es
\TiempoAgotamientoMedio{} h, con IC 95\%
[\TiempoAgotamientoICInf{}, \TiempoAgotamientoICSup{}] h.

\begin{figure}[H]
\centering
\includegraphics[width=.78\linewidth]{report_assets/figuras/trayectorias_criticas.pdf}
\caption{Inventario final medio e IC 95\% de las combinaciones críticas en los escenarios base y post-intercambio.}
\label{fig:trayectorias}
\end{figure}

\begin{table}[H]
\centering
\caption{Primer agotamiento crítico en zonas abastecibles.}
\label{tab:agotamiento}
\small
\input{report_assets/tablas/tabla_agotamiento.tex}
\end{table}

El balance completo conserva una fila por bloque, zona y suministro, con
inventario inicial, reposición, demanda, consumo, inventario final, déficit y
superávit; su versión final se exporta en el paquete de activos. La demanda media total
cambia de \DemandaBase{} a \DemandaPost{} unidades y el déficit medio total de
\DeficitBase{} a \DeficitPost{}. La comparación incluye IC 95\% calculados
sobre realizaciones, no sobre filas mezcladas.

\begin{table}[H]
\centering
\caption{Comparación agregada antes y después del intercambio.}
\label{tab:comparacion}
\small
\input{report_assets/tablas/tabla_comparacion.tex}
\end{table}

La Figura~\ref{fig:deficit} separa la brecha estructural de Z1 ---sin centro ni
ruta de entrada--- de los déficits operativos. La cobertura estructural final es
\CoberturaFinal{} de las zonas del contrato.

\begin{figure}[H]
\centering
\includegraphics[width=.68\linewidth]{report_assets/figuras/deficit_zonal_post.pdf}
\caption{Déficit medio acumulado por zona en el escenario post-intercambio.}
\label{fig:deficit}
\end{figure}

\subsection{Convoy, distribución y cuellos de botella}

En el escenario base, el último momento seguro para activar el convoy deja un
margen de \MargenConvoyBase{} h. Para la nueva zona crítica post-intercambio el
margen factible es \MargenConvoy{} h: no existe un convoy externo que llegue
antes del agotamiento con la red disponible. Esta diferencia evita interpretar
una mejora temporal promedio como factibilidad logística automática.

La ruta prioritaria es \RutaPrioritaria{}, con
\ViajesRutaPrioritaria{} viajes por bloque. La asignación usa
\CombustibleUtilizado{} galones y nunca excede la reserva por tipo de vehículo.
La frecuencia es una recomendación heurística, pues faltan equivalencias entre
metros cúbicos y las unidades de suministros.

\begin{table}[H]
\centering
\caption{Primeras rutas de la recomendación de distribución.}
\label{tab:plan}
\scriptsize
\input{report_assets/tablas/tabla_plan.tex}
\end{table}

\begin{table}[H]
\centering
\caption{Principales cuellos de botella trazados a los datos.}
\label{tab:cuellos}
\small
\input{report_assets/tablas/tabla_cuellos.tex}
\end{table}

\textbf{Respuesta 1 del Grupo 3.} Antes del intercambio, el primer suministro
crítico es medicamentos en \ZonaCriticaBase{}; el convoy dispone de
\MargenConvoyBase{} h de margen. Después del intercambio, entre zonas
abastecibles, el primer crítico es \SuministroCritico{} en
\ZonaCriticaPost{}, en el bloque asociado a \TiempoAgotamientoMedio{} h, y no
hay convoy externo factible a tiempo.

\textbf{Respuesta 2 del Grupo 3.} La zona crítica abastecible sí cambia de
\ZonaCriticaBase{} a \ZonaCriticaPost{}. La atribución por escenarios aislados
señala como factor dominante la demanda formal del Grupo 1. La accesibilidad del
Grupo 4 modifica tiempos y factibilidad del convoy, pero el motor de balance no
usa rutas para alterar inventario; por tanto no se le atribuye un cambio de
déficit que el código no modela.

\section{Incorporación del intercambio}

\textbf{Grupo 1:} se recibió población por Z1--Z5 y bloque. La estructura ancha
se transformó a bloque--zona y sustituyó la estimación proporcional a capacidad
de almacenamiento. Con las mismas tasas del Excel se recalculó
$D_{zst}=P_{zt}c_s$. Esto reveló demanda real en Z1 y cambió la combinación
crítica abastecible.

\textbf{Grupo 4:} se recibió vehículo transitable, tiempo mediano desde Z4 e
índice de aislamiento por zona y bloque. Los tiempos solo sustituyen rutas con
origen en la Bodega Central de Z4; el código liviano impide asignar camión
grande. No se alteran capacidades porque el reporte no las proporciona, ni se
crean rutas para Z1 o para destinos cuya zona sigue sin identificar.

\begin{center}
\fbox{\parbox{.92\linewidth}{
\textbf{Antes} (población estimada y tiempos del Excel)
$\rightarrow$ \textbf{información recibida} (población formal y accesibilidad)
$\rightarrow$ \textbf{cambio del modelo} (demanda y restricciones de tránsito)
$\rightarrow$ \textbf{después} (\ZonaCriticaPost{}/\SuministroCritico{},
con brecha estructural en Z1).}}
\end{center}

La trazabilidad completa ---campo recibido, transformación, regla afectada y
supuesto confirmado o contradicho--- se exporta con el Excel final. Esta cadena
separa el dato externo de las decisiones de equivalencia del equipo.

\section{Limitaciones y propuestas de mejora}

\begin{tabularx}{\linewidth}{@{}p{.19\linewidth}X X@{}}
\toprule
\textbf{Limitación} & \textbf{Impacto y posible sesgo} & \textbf{Mejora y dato necesario}\\
\midrule
Capacidad en m$^3$ & No puede restringirse rigurosamente agua, raciones y kits; la frecuencia puede sobrestimar cobertura. & Incorporar volumen y compatibilidad de carga por unidad de cada suministro.\\
Reposición sin centro & Asignarla a la Bodega Central concentra disponibilidad y puede sesgar el agotamiento zonal. & Registrar centro receptor y tiempo de traslado para cada reposición.\\
Bloqueo parcial & No existe porcentaje residual; conservar capacidad puede ser optimista. & Obtener capacidad y velocidad por ruta y bloque, vinculadas inequívocamente al mapa del Grupo 4.\\
Albergue Estadio & Sin zona no puede competir por prioridad ni recibir accesibilidad zonal. & Entregar coordenadas o equivalencia oficial destino--zona.\\
Información perfecta & La prioridad oficial supone observación inmediata y puede reaccionar tarde en operación. & Integrar telemetría y una política de despacho que retroalimente inventarios.\\
\bottomrule
\end{tabularx}

La sensibilidad de seis horas conserva el déficit físico, porque el plan no
retroalimenta el balance, pero deja \DeficitNoObservadoRetraso{} unidades sin
observar al cierre. Cambian \RutasCambioPrioridad{} prioridades; la zona más
afectada por información tardía es \ZonaAfectadaRetraso{}. Presentar ese déficit
como ``adicional causado'' sería ficticio: cuantificarlo exige que los despachos
modifiquen inventarios y cobertura dentro del motor.

\begin{table}[H]
\centering
\caption{Sensibilidad a seis horas de retraso informativo.}
\label{tab:sensibilidad}
\small
\input{report_assets/tablas/tabla_sensibilidad.tex}
\end{table}

\end{document}
